% !TEX program = lualatex
\documentclass[lualatex,ja=standard]{bxjsarticle}

\usepackage{amsmath}
\usepackage[hidelinks]{hyperref}
\usepackage{tikz}

\title{相互参照のサンプル}
\author{LaTeX サンプル}
\date{\today}

\begin{document}

\maketitle
\tableofcontents

% -------------------------------------------------------------------
% 参照先の定義（ここでラベルを貼る）
% -------------------------------------------------------------------
\section{数式}
\label{sec:math}

\begin{equation}
  e^{i\pi} + 1 = 0
  \label{eq:euler}
\end{equation}

\begin{equation}
  \int_0^1 x^2 \, dx = \frac{1}{3}
  \label{eq:integral}
\end{equation}

\section{図}
\label{sec:fig}

\begin{figure}[h]
  \centering
  \begin{tikzpicture}
    \draw[thick, fill=green!30] (0,0) -- (2,0) -- (1,1.7) -- cycle;
    \node at (1, -0.3) {三角形};
  \end{tikzpicture}
  \caption{サンプルの三角形}
  \label{fig:triangle}
\end{figure}

\section{表}
\label{sec:tab}

\begin{table}[h]
  \centering
  \caption{サンプルの表}
  \label{tab:sample}
  \begin{tabular}{|l|r|}
    \hline
    項目 & 値 \\
    \hline
    A    & 10 \\
    B    & 20 \\
    \hline
  \end{tabular}
\end{table}

% -------------------------------------------------------------------
% 参照のまとめ（ここで一気に参照する）
% -------------------------------------------------------------------
\section{相互参照}
\label{sec:ref}

上で定義したラベルをここでまとめて参照します。

\subsection{\texttt{\textbackslash ref} と \texttt{\textbackslash pageref}}

\texttt{\textbackslash ref\{キー\}} で番号を、\texttt{\textbackslash pageref\{キー\}} でページ番号を参照します。

\begin{itemize}
  \item セクション参照: \ref{sec:math} 節（\pageref{sec:math} ページ）
  \item 数式参照: 式 \ref{eq:euler}
  \item 図参照: 図 \ref{fig:triangle}
  \item 表参照: 表 \ref{tab:sample}
\end{itemize}

\subsection{\texttt{\textbackslash eqref}（数式専用）}

\texttt{\textbackslash eqref} は数式番号を括弧付きで参照します。

\begin{itemize}
  \item 式~\eqref{eq:euler} はオイラーの等式
  \item 式~\eqref{eq:integral} は定積分の例
\end{itemize}

\subsection{\texttt{\textbackslash autoref}（種別を自動付与）}

\texttt{\textbackslash autoref} は「節」「図」「表」などのプレフィックスを自動で付加します。

\begin{itemize}
  \item \autoref{sec:math}
  \item \autoref{fig:triangle}
  \item \autoref{tab:sample}
  \item \autoref{eq:euler}
\end{itemize}

\subsection{\texttt{\textbackslash nameref}（セクション名で参照）}

\texttt{\textbackslash nameref} はセクション番号の代わりにセクション名を表示します。

\begin{itemize}
  \item \nameref{sec:math}（\autoref{sec:math}）
  \item \nameref{sec:fig}（\autoref{sec:fig}）
  \item \nameref{sec:tab}（\autoref{sec:tab}）
\end{itemize}

\end{document}
